\cleardoublepage

\section{绪论}

\subsection{研究背景与意义}

近年来，大规模预训练模型的发展显著推动了视觉理解与跨模态推理技术的进步。以 Transformer 为代表的深度学习架构通过自注意力机制实现了对长序列信息的高效建模\cite{transformer2017}，并逐渐成为自然语言处理与计算机视觉领域的主流技术路线。在视觉领域，Vision Transformer（ViT）通过将图像划分为 patch token 并进行序列化建模，验证了 Transformer 在视觉任务中的有效性\cite{vit2021}。随后，CLIP 模型利用大规模图文对比学习实现视觉语义与文本语义空间的统一表示，为多模态大模型的发展奠定了重要基础\cite{clip2021}。

在此基础上，多模态大模型逐渐从静态图像理解扩展至统一视觉语言推理。Flamingo 通过跨模态注意力结构实现视觉与文本的联合生成\cite{flamingo2022}；BLIP-2 提出利用轻量级 Q-Former 模块建立视觉编码器与大型语言模型之间的桥接结构，在降低训练成本的同时保持较强的视觉推理能力\cite{blip22023}；LLaVA 与 MiniGPT-4 等工作则通过视觉指令微调进一步提升模型的复杂推理与交互能力\cite{llava2023,minigpt42023}。目前，“视觉编码器—视觉投影器—大型语言模型”已逐渐成为多模态大模型的主流结构范式。

随着研究逐渐向视频场景扩展，视频多模态大模型开始具备视频问答、时空推理以及动态场景理解能力。Video-LLaMA 首次实现音频、视频与文本三模态联合建模\cite{videollama2023}；Video-LLaVA 探索统一视觉表示方式以增强视频场景中的视觉泛化能力\cite{videollava2023}；VideoLLaMA2 则通过引入 STC（Spatio-Temporal Convolution）Connector 强化视频长时序关系建模能力\cite{videollama22024}。上述研究推动了视频多模态模型从简单帧级特征拼接向专用时空结构设计的演进。

然而，视频多模态模型在获得更强理解能力的同时，也面临显著的推理效率问题。视频输入通常包含大量连续帧，视觉 token 数量会随着视频长度和分辨率快速增长，导致模型推理计算复杂度显著增加。对于基于 Transformer 的架构，自注意力模块的计算复杂度通常与 token 数量呈平方关系，因此大量视觉 token 会带来较高的显存占用与推理延迟\cite{timesformer2021,vivit2021}。此外，在视频多模态模型中，视觉 token 还需进一步参与跨模态融合与语言生成过程，从而进一步放大计算与存储压力。

当前视频场景中普遍存在明显的视觉冗余与时空冗余。例如，相邻视频帧之间通常具有较高的语义相似性，大量 patch token 在推理阶段会产生重复计算，但现有模型仍需对其进行完整编码与注意力计算。这种冗余不仅影响推理效率，也限制了视频大模型在实时场景中的部署能力。因此，如何在尽量保持视频理解性能的前提下减少视觉 token 数量、降低推理复杂度，已成为视频多模态大模型研究中的重要问题。

针对上述问题，近年来研究开始探索视觉 token 压缩与动态计算优化方法。Token Merging（ToMe）通过合并相似视觉 token 降低 Transformer 输入长度\cite{tome2023}；DynamicViT 利用动态 token 剪枝策略保留重要视觉区域\cite{dynamicvit2021}；PruneVid 则尝试将视觉 token 剪枝扩展至视频大模型场景\cite{prunevid2024}。此外，TokenPacker 提出一种 coarse-to-fine 的视觉 token 压缩方案，通过多层细节注入机制，在压缩视觉 token 的同时尽可能保留局部细粒度语义信息\cite{tokenpacker2024}。

基于上述背景，本文选择以 VideoLLaMA2-7B-16F 为基础框架，将 TokenPacker 模块迁移至视频大模型结构中，在视觉编码器与 STC Connector 之间实现逐帧空间视觉 token 压缩。本文在保持原有时空建模结构与大型语言模型主体不变的前提下，通过引入 TokenPacker 对视频帧中的空间视觉 token 进行压缩，以减少后续时空融合与语言解码阶段的计算负担，从而提升视频多模态模型的推理效率。

\subsection{国内外研究现状}

\subsubsection{视频多模态大模型研究现状}

随着多模态大模型的发展，研究重点逐渐从静态图像理解扩展至视频理解与视频推理任务。相较于图像场景，视频不仅包含空间语义信息，还具有复杂的时间动态关系，因此模型需要同时具备空间建模与时序建模能力。

早期多模态模型主要围绕视觉与文本统一表示展开研究。CLIP 通过大规模图文对比学习实现视觉与语言特征的统一对齐\cite{clip2021}；Flamingo 引入跨模态注意力机制，实现视觉与文本的联合生成\cite{flamingo2022}；BLIP-2 则通过 Q-Former 构建视觉编码器与大型语言模型之间的桥接结构\cite{blip22023}。随后，LLaVA 与 MiniGPT-4 等模型进一步利用视觉指令微调策略增强复杂视觉推理能力\cite{llava2023,minigpt42023}。

随着研究逐渐向视频场景扩展，Video-LLaMA 首次实现音频、视频与文本三模态联合建模\cite{videollama2023}；Video-LLaVA 通过统一视觉表示方式提高视频场景中的视觉泛化能力\cite{videollava2023}；VideoLLaMA2 则通过引入 STC Connector 增强长时序视频建模能力\cite{videollama22024}。此外，InternVideo 等工作通过大规模视频预训练进一步提升视频特征表达能力\cite{internvideo2023}。

尽管现有视频多模态模型在视频问答、时空推理与长视频理解任务中取得较好效果，但随着视频帧数与分辨率增加，视觉 token 数量快速增长，导致推理延迟与显存占用显著增加。因此，视频多模态模型的高效推理问题逐渐成为当前研究的重要方向。

\subsubsection{视觉Token压缩研究现状}

为降低 Transformer 模型的计算复杂度，研究者提出了多种视觉 token 压缩与动态剪枝方法。Token Merging（ToMe）通过合并相似视觉 token 减少序列长度，从而降低 Transformer 的注意力计算量\cite{tome2023}；DynamicViT 则利用动态 token 剪枝策略，在推理过程中保留重要视觉区域\cite{dynamicvit2021}。

近年来，视觉 token 压缩研究逐渐扩展至多模态大模型场景。FastV 通过动态视觉 token 压缩降低多模态模型中的视觉计算冗余；PruneVid 则尝试在视频大模型场景中进行视觉 token 剪枝\cite{prunevid2024}。此外，一些方法通过 pooling 或卷积下采样减少视觉 token 数量，以降低推理复杂度。

TokenPacker 提出了一种 coarse-to-fine 的视觉 token 压缩方案\cite{tokenpacker2024}。该方法首先构建低分辨率粗粒度查询，再利用多层高分辨率区域特征进行细节注入，从而在压缩视觉 token 的同时尽可能保留局部语义信息。相比于简单 pooling 或直接 token 下采样方法，TokenPacker 在图像多模态场景中表现出较好的效率与性能平衡能力。

\subsubsection{视频场景高效推理研究现状}

围绕视频大模型高效推理问题，当前研究主要从帧采样、时空建模优化、稀疏注意力以及长视频推理机制等方向展开。

在视频建模方面，TimeSformer 将视频建模中的空间注意力与时间注意力分离，从而降低视频 Transformer 的计算复杂度\cite{timesformer2021}；ViViT 探索纯 Transformer 架构在视频理解中的应用\cite{vivit2021}。此外，StreamingLLM 与 LongNet 等工作分别从流式推理与超长上下文建模角度提升长序列推理效率\cite{streamingllm2023,longnet2023}。

在长视频理解方向，LongVLM、LLaMA-VID 等研究尝试通过减少视觉 token 数量与优化视觉表示结构缓解长视频场景中的 token explosion 问题\cite{longvlm2024,llamavid2024}。此外，部分研究还探索了动态帧采样、关键帧筛选与时空联合压缩等方法，以减少视频推理中的冗余计算。

尽管上述研究在提升视频多模态模型推理效率方面取得了一定进展，但仍存在一些不足。首先，现有方法大多聚焦于帧级压缩或注意力结构优化，对于视频中大量空间冗余 token 的利用仍不充分；其次，部分方法在高压缩倍率场景下容易出现明显性能下降；此外，目前关于“逐帧空间压缩+后续时空融合”的系统性研究仍较少。因此，探索一种能够适应视频时空结构特点的视觉 token 压缩方案，对于提升视频多模态模型推理效率具有重要意义。

\subsection{本文研究内容}

针对视频多模态大模型在推理过程中存在视觉 token 数量庞大、推理延迟高以及显存占用大的问题，本文围绕 VideoLLaMA2-7B-16F 模型开展视觉 token 压缩研究，重点探索 TokenPacker 从图像场景向视频场景迁移的可行性与有效性。

本文首先分析 VideoLLaMA2 的视觉分支结构与 STC Connector 的工作机制，梳理视觉 token 在视频模型中的流转过程，并分析视频场景中的空间冗余问题。在此基础上，本文将 TokenPacker 模块迁移至 VideoLLaMA2 中，在视觉编码器与 STC Connector 之间构建逐帧空间 token 压缩方案，实现对视频帧中空间视觉 token 的压缩与重组织。

随后，本文针对直接迁移导致的视频理解性能下降问题，进一步设计 Token 校准与轻量微调策略，通过对齐压缩后视觉 token 分布与微调压缩模块参数，缓解视觉特征分布偏移问题，提高模型在视频任务中的性能表现。

最后，本文通过压缩效率实验、模型性能对比实验、不同压缩倍率下的鲁棒性实验以及泛化能力实验，对所提出方法进行验证，并分析不同压缩倍率下模型性能变化规律以及 TokenPacker 在高压缩率场景中的适用性。

\subsection{本文创新点}

本文围绕视频多模态大模型高效推理问题开展研究，主要创新点如下：

（1）提出 TokenPacker 在视频大模型中的迁移方案。在保持 VideoLLaMA2 原有 STC Connector 与大型语言模型主体结构不变的前提下，将 TokenPacker 引入视觉编码器与 STC Connector 之间，实现视频场景中的视觉 token 压缩。

（2）设计逐帧空间视觉 token 压缩策略。本文采用“先逐帧空间压缩，再进行时空融合”的方案，在不破坏原有时空建模结构的条件下减少视觉 token 数量，从而降低后续时空聚合与语言解码阶段的计算开销。

（3）构建 Token 校准与轻量微调流程。针对 TokenPacker 直接迁移导致的视频理解性能下降问题，本文设计 Token 校准与轻量微调策略，以缓解视觉特征分布偏移问题并恢复模型性能。

（4）验证高压缩倍率场景下的鲁棒性优势。实验结果表明，相较于 Mean Pooling 等简单压缩方法，本文方法在高压缩倍率场景下准确率下降更平缓，表现出更好的压缩鲁棒性。

\subsection{论文结构安排}

本文共分为五章，各章节内容安排如下：

第一章为绪论，主要介绍课题研究背景、研究意义、国内外研究现状以及本文的研究内容与创新点。

第二章为相关理论与关键技术，主要介绍 Transformer 与注意力机制、CLIP 视觉编码器、VideoLLaMA2 模型结构以及 TokenPacker 原理等内容，为后续方法设计提供理论基础。

第三章为方案设计，重点介绍本文提出的整体架构设计、视频视觉特征构建、TokenPacker 视频迁移方案、Token 校准与轻量微调策略以及复杂度分析。

第四章为实验与结果分析，主要通过压缩效率实验、模型性能对比实验、不同压缩倍率下的鲁棒性实验以及泛化能力实验，对本文方法进行验证与分析。

第五章为总结与展望，对本文工作进行总结，并分析当前方法存在的不足以及未来研究方向。

